% =============================================================================
% Results sec 5.7 -- Mechanism and robustness checks
% Compressed version. Target ~390 words.
% =============================================================================

\section{Mechanism and robustness checks}
\label{sec:mechanism-robustness}


\subsection*{\S\,5.7.1 Update-norm activity}

FedProx has two distinct mechanisms: a per-client proximal regulariser
that operates during local training, and a protocol-level $\gamma$-inexact
policy that operates at aggregation. The centrepiece
\S\,\ref{sec:li-straggler} result attributes the straggler-regime win to
the latter. The engineered update-norm reduction reported here reflects
the former: on the engineered partition ($n=10$), mean per-round
update-norm magnitude $\overline{\lVert \Delta w \rVert}$ is $1.523$ for
FedAvg vs $1.029$ for FedProx ($\mu = 0.01$) --- a $\approx 32\,\%$
reduction, consistent with the proximal term's drift-control intent. This
is a per-protocol observation of drift control under symmetric conditions,
distinct from the \S\,\ref{sec:li-straggler} aggregation-level mechanism
that carries the straggler-regime win.

Across 78 L4 runs pooled over eight algorithm-by-protocol cells (baselines,
Li 2020 four-condition, perfect-storm, asymmetric $\mu$, asymmetric LR,
weighted-CE), dominant-client mean update-norm magnitude correlates with
rare-class F1 at Spearman $r = +0.634$. The relative-influence ratio of the
rare-class-holding client (its update-norm $\div$ Client 0's) correlates
near zero across protocols, $r = -0.045$: it is not by itself predictive
when the comparison spans heterogeneous protocols.

Within the asymmetric-LR experiment alone ($n=27$ cells), the relative-
influence ratio does track rare-class F1 closely: $r = +0.80$ for FedAvg,
$r = +0.87$ for FedProx (cf.\ \S\,\ref{sec:lr-asymmetry}). Within-protocol
ordering does not extrapolate across the eight-cell L4 protocol family. We
report these correlations as exploratory mechanism evidence --- consistent
with a FedProx drift-control and a FedNova contribution-preservation
interpretation --- not as causal proof. The causal claim rests on
\S\,\ref{sec:li-straggler}.


\subsection*{\S\,5.7.2 Protocol-flip audit}

A protocol-flip audit across six multi-seed L4 cells under three
measurement choices (best-validation checkpoint, final-round, last-$K$-round
mean): large-effect cells (perfect-storm, Li 2020 asymmetric,
asymmetric-$\mu$) flip $0/9$ seed-level comparisons; small-effect cells
(node-pinned L4 s42, extended-rounds L3 s42, two-client 90/10 s456) flip
$3/9$. Small effects in symmetric protocols are reporting-protocol-fragile;
large asymmetric-protocol effects are not.

The chapter therefore weights conclusions by effect magnitude: the
\S\,\ref{sec:li-straggler} centrepiece ($+0.115$ macro-F1) and the
\S\,\ref{sec:lr-asymmetry} 50:1 FedNova--FedAvg gap ($+0.131$) are robust
to checkpoint choice; the sub-$1$\,pp differences in
\S\,\ref{sec:statistical-heterogeneity} are not. \S\,5.8 labels each
finding's strength accordingly.

The $B$-dissimilarity proxy, training-instability metric, asymmetric-$\mu$
Yao test, early-warning collapse detector, and small-hospital case study
are reported in Appendix~E as exploratory; none is load-bearing for any
\S\S\,5.2--5.8 conclusion.

These checks sketch a regime map for the three optimisers, summarised next.


% =============================================================================
% Figure 5.8 placeholder -- update-norm mechanism
% =============================================================================

\begin{figure}[h]
\centering
\includegraphics[width=0.85\textwidth]{F_update_norms.pdf}
\caption{Update-norm mechanism on the engineered partition. Mean per-round
$\lVert \Delta w \rVert$ trajectories per client for FedAvg vs FedProx
($n=10$ seeds); FedProx tracks lower throughout, consistent with the
proximal term's drift-control intent. Cross-protocol scatter panel
(78 L4 runs across 8 algorithm $\times$ protocol cells): Spearman
$r = +0.634$ between dominant-client mean update-norm and rare-class F1;
relative-influence ratio $r = -0.045$ across protocols.}
\label{fig:update-norms}
\end{figure}
